\section{Introduction}\label{sec:intro}

\subsection{The problem}

Can a deterministic multitape Turing machine transpose a $d\times d$ matrix with $b(d)=\max(1,\clog d)$-bit entries, written
row by row, in time linear in the input length $L=d^2b(d)$? This is the statement~(T) of the first paper of this
series~\cite{FirstPaper2026}. By a reduction of Harvey and van der Hoeven~\cite[Cor.~6.9]{HvdH2025}, a linear-time
multiplication machine would give a linear-time transposer, so a proof of $\neg$(T) would give the first superlinear lower
bound for off-line integer multiplication on multitape Turing machines. The question is open, and this paper does not decide it.

For oblivious machines the fourth paper of the series~\cite{FourthPaper2026} reduced $\neg$(T) to two statements of different
kinds (its Theorem~D). The first is \emph{depth reduction}: the time--space dependency DAG of every oblivious linear-time run
can be cut to bounded depth by revealing $\eps L$ bits. At the level of head walks this is the \emph{trajectory cut problem}.
A \emph{walk} is a sequence $x(0),x(1),\dots$ of integers with $\abs{x(t+1)-x(t)}\le1$, the positions of one head; a
\emph{long hop} is a pair $p<u$ of consecutive visits to one cell with $u\ge p+2$. The \emph{walk DAG} of $k$ walks of length~$T$
has the vertices $0,\dots,T-1$, the state edges $t\to t-1$, and an edge $u\to p$ for every long hop of every walk. $\TCP(k)$
says: for every $\eps>0$ there is a $D$ such that, for every $T$ and every $k$ walks of length~$T$, deleting some $\eps T$ edges
leaves no path with more than $D$ hop edges. The second statement, the \emph{crux}, is about codes: every index code for the
side information of the remaining bounded-depth cones needs more than $\eps L$ bits.

Depth reduction has been settled step by step. The fourth paper proved TCP for home-based work heads, whose legs all begin or
end near fixed homes, for every coupling of the heads (its Theorem~B). The eighth paper~\cite{EighthPaper2026} proved it for
walks of bounded nesting load (its Theorem~F), and for one arbitrary walk together with walks of bounded load at boundedly many
scales (its Theorem~M). It also showed that TCP is, up to a factor~$2$ in~$k$, depth reduction for \emph{stack DAGs}
(\cite[Proposition~4.5]{EighthPaper2026}; \cref{prop:stacksred} here). An $s$-stack DAG on $[0,N)$ has the spine edges
$t\to t-1$ and arcs $u\to p$ with $u\ge p+2$, partitioned into $s$ families in which no two arcs cross and each vertex is the
left end of at most one arc and the right end of at most one. $\SR(s)$ says that every family of $s$-stack DAGs is
depth-reducible in the above sense. Two stacks are planar, so $\SR(2)$ holds; $\SR(3)$ was the first open
case~\cite[Problem~10.2]{EighthPaper2026}.

\subsection{The expected answer, and the actual one}

The series expected $\TCP(2)$ to fail. Its sharpest candidate for a depth-robust family of walks was the \emph{deep-nesting
pair} of the fourth paper: a walk whose strokes are nested at $m+1$ levels, together with a copy shifted in time. Conjecture~A$'$
of the fourth paper and its weakening, Conjecture~A$''$ of the eighth, would make these pairs depth-robust as $m$ grows, with depth
$2^{\Omega(m)}$ at fixed cost~\cite[Corollary~5.26]{EighthPaper2026}; lower bounds for several restricted classes of labels
supported this expectation~\cite{FourthPaper2026,EighthPaper2026,DirectedReciprocity2026,MergeDepth2026}. The actual answer is the
opposite: $\TCP(k)$ holds for every~$k$.

\subsection{Results}

\begin{itemize}[leftmargin=1.5em]
\item \textbf{Theorem~T} (\cref{thm:T}). For every $k\ge1$ and $\eps>0$ there is a $D(k,\eps)$ such that the walk DAG of any $k$
walks of any length~$T$, with or without pauses, has a cut of at most $\eps T$ edges that leaves no path with more than
$D(k,\eps)$ hop edges. So $\TCP(k)$ holds for every~$k$, and by the reduction of the eighth paper \textbf{$\SR(s)$ holds for every
$s$}: bounded-stack DAGs with a spine reduce to constant depth. The bound is explicit, $D(k,\eps)=K^P-1$ with the parameters of
\cref{def:params}, and $\log\log D(k,\eps)=O(k^2\log(k/\eps))$ for $\eps\le\frac12$ (\cref{sec:open}).
\item \textbf{Corollary~N} (\cref{cor:dnp}). The deep-nesting pairs are cut to depth $D(2,\eps)$, independently of the nesting
depth~$m$ and of the shift. Hence Conjectures~A$'$ and~A$''$ are false: on a cyclic cover, the label of Theorem~T is cheap and
doubly serves almost every level with a number of values independent of~$m$. Both statements are Lean-verified.
\item \textbf{Corollary~O} (\cref{cor:oblivious}). Every oblivious machine with $k$ work tapes running in time $C(L+1)$ has, for
every $\eps>0$, a genie of cost at most $\eps L$ whose full depth is bounded in terms of $k$, $C$, the value width and~$\eps$. So
the crux alone excludes oblivious linear-time transposers, whatever their head motion, and for machines with a linear data path
the linear crux alone does. These statements are Lean-verified in the frozen machine model of the first paper. In particular, the native form of Question~9.1 of \emph{Directed Recovery
Structures}~\cite{DirectedRecovery2026}, the half bound for all codes on the cone relations of linear-time machines, would exclude
them (\cref{rem:native}). On the side of constructions, an oblivious linear-time transposer must carry codes of rate close to~$1$
on the cone relations of its own schedule at bounded depth: the coding side is the whole question (\cref{sec:cons-construction}).
\item \textbf{Corollary~HV} (\cref{cor:HV}). The same genies give the ball systems of oblivious linear-time transposers index
codes of at most $\eps L$ bits at a bounded radius, and, for machines with a linear data path, minrank at most $\eps L$. So
Conjecture~HV$^*$ of the eleventh paper, code-level depth reduction, holds, and its Problem~9.1 has an affirmative
answer~\cite{MergeDepth2026}. This is Lean-verified as well.
\item \textbf{Corollary~L} (\cref{cor:L}). Given the native half bound for the cone relations of frozen traces, no machine
transposes in linear time with a trace of entropy at most $\eta L$, $\eta<\frac12$, however often it enters its cells: this is
the eighth paper's Corollary~F$''$ without its bounded-entry hypothesis. The eleventh paper's Theorem~SE loses its bounded-entry
hypothesis in the same way. Both are Lean-verified, with the native bound as an explicit hypothesis.
\end{itemize}

\subsection{How the proof works}

The proof has three layers.

\emph{Cuts are labels.} A labelling $\Psi:[0,T)\to\{0,\dots,q-1\}$ defines a cut: the state edges where $\Psi$ decreases and the
hop edges $u\to p$ with $\Psi(p)\ge\Psi(u)$. Every path avoiding the cut increases $\Psi$ at each hop edge, so its depth is less than
$q$ (\cite{EighthPaper2026}; \cref{prop:labelcut} here). We build a random label with $q=K^P$ values and expected cut at most
$\eps T$; some outcome gives the cut (\cref{prop:avg}).

\emph{Ladders, one walk at a time} (\cref{sec:ladders}). The label gives each hop a \emph{deciding group}, a coarse scale, and
pays for the hops whose two ends see different scales. To make this rare we measure scales from homes adapted to the walk. Every
upper hop has a \emph{top}, the first time inside it where the walk is highest; the hops with a common top are nested, one per
cell, and form the \emph{ladder} of the top. Their ends are the last visits to each cell before the top and the first visits after
it. The ladders partition the upper hops. The outermost hop of a ladder, the \emph{canonical excursion} of its top, starts and
ends at the \emph{base}~$y$ of the ladder; the walk between consecutive rungs forms a \emph{hanging structure} on the climb and a
\emph{descent structure} on the descent, and the canonical excursion of every other top that it contains lies in one of these
structures. Measuring the distance from the base of each excursion, suspending the structures as \emph{frames} and restarting
the descent at the levels $y+2^e$, makes every hop cheap on average over a random grid: \emph{Theorem~A} (\cref{thm:A}) bounds
the expected number of hops lost per walk by $6$ times the number of hops divided by the grid width~$h_0$.

\emph{Displays with break bits, all walks together} (\cref{sec:label,sec:displays,sec:accounting}). Frames create an interaction
between walks: while a frame runs, the segment that it suspends keeps its band frozen, and the band reappears when the segment
resumes. The label therefore shows, besides the current band group of each walk, a \emph{display} of suspended frames, and
digits that count events at each shown scale.
\emph{Lemma~V} (\cref{lem:V}) shows that a hop whose deciding group is shown at both ends increases the label unless the display
is truncated, the shown set above the hop's group changes, or a digit wraps. The second case is the heart of the matter.
\emph{View invariance} (\cref{prop:W}) shows that the part of the display above a group depends only on a \emph{view} that changes
only at standard events, which are paid for; the display is defined by random \emph{break bits}, and the hops exposed despite
view invariance are those with a broken candidate, at most a $2n/Q$ fraction (\cref{prop:exposure}). Exact displays would be
unbounded (\cref{rem:exact}); with break bits, \emph{Lemma~D} (\cref{lem:D}) bounds the size of the display independently of the
depth of the stacks.

\subsection{Formal status and how to read the paper}

Theorem~T, every numbered result of \cref{sec:prelim,sec:ladders,sec:label,sec:displays,sec:accounting,sec:proof}, and
Corollaries~N, O, HV and~L with their definitions are formalized in Lean~4 with Mathlib, in the library
\leanfile{BoundedStackDepth} next to this paper, which builds on the eighth paper's library and, for the corollaries, on the
eleventh paper's library~\cite{MergeDepth2026}, which formalizes the deep-nesting walk and its cyclic covers, and traces,
dependency DAGs and index codes in the frozen machine model of the first paper. Each result is followed by a status line naming its Lean declarations and file, and by an outline of the
formal proof; the Lean statement is authoritative. The formal main theorem is \lean{tcpAt_direct}: for $k$ walks with
pauses, cuts of at most $\eps T$ edges and depth at most $K^P-1$. Its corollaries are \lean{tcp} and \lean{main} ($\TCP(k)$ for
every~$k$, the statement \lean{MainStatement} of \leanfile{Statement.lean}) and \lean{sr} ($\SR(s)$ for every~$s$, through the
eighth paper's reduction). All of them depend only on the axioms \texttt{propext}, \texttt{Classical.choice} and \texttt{Quot.sound} (\cref{app:verify}). The
results of the eighth paper that the proof uses are Lean-verified there and are marked as cited.

Not formalized are the remarks of \cref{sec:ingredients,sec:consequences}. They explain why simpler constructions fail, or
discuss the consequences; they are labelled as proof sketches, heuristics or numerical observations, and no proof uses them.

\subsection{Related work}

\emph{Depth-robust graphs.} Erd\H{o}s, Graham and Szemer\'edi~\cite{EGS1975} constructed sparse DAGs in which every large set of
vertices leaves a long path; Valiant~\cite{Valiant1977} used depth reduction as a route to lower bounds for circuits;
Schnitger~\cite{Schnitger1983} constructed depth-robust graphs of constant degree; and depth-robust graphs are a basic tool for
memory-hard functions~\cite{ABP2017}. Theorem~T says that no family of bounded-stack DAGs with a spine is depth-robust, when
only arcs are counted along paths.

\emph{Time and space.} The computation graphs of multitape machines were studied through pebbling by Hopcroft, Paul and
Valiant~\cite{HPV1977}, through edge-deletion depth reduction by Paul and Reischuk~\cite{PaulReischuk1980}, and through
segregators, whose removal leaves every vertex with few ancestors, by Paul, Pippenger, Szemer\'edi and Trotter~\cite{PPST1983}
and by Santhanam~\cite{Santhanam2001}. Theorem~T is a statement of another kind about the walk DAGs of the work heads: after
the removal of $\eps T$ edges, every path uses a bounded number of long hops, although it may use many state edges.

\emph{Layouts.} Galil, Kannan and Szemer\'edi~\cite{GKS1989separators,GKS1989pushdown} studied separators of $k$-page graphs
and showed that $3$-page graphs can lack small separators. Some bounded-degree expanders have $3$-page and $2$-queue
layouts~\cite{DSW2016}, and monotone expanders exist~\cite{DvirWigderson2010,BourgainYehudayoff2013}. Expansion is not
depth-robustness in our sense: the spine is free, and only arcs are counted.

\emph{Long paths.} The decomposition of a tree into long paths is a standard data-structure tool, used for example for level
ancestors~\cite{BenderFarachColton2004}, where the paths extended upwards are also called ladders. Our ladders are such a
decomposition, without the extension, of the forest of nested upper hops; its use to bound label costs appears to be new.

\subsection{Organization}

\Cref{sec:prelim} fixes walks, hops, cuts, prominence and the random grid. \Cref{sec:ladders} proves Theorem~A, one walk at a
time. \Cref{sec:label} defines the label and proves Lemma~V. \Cref{sec:displays} defines displays with break bits and proves
Lemma~D, view invariance and the exposure bound. \Cref{sec:accounting} adds up the costs, and \cref{sec:proof} chooses the
parameters and proves Theorem~T. \Cref{sec:ingredients} explains why each ingredient is needed, \cref{sec:consequences} draws the
consequences for tape machines, and \cref{sec:open} lists open problems. \Cref{app:index} indexes the Lean declarations, and
\cref{app:verify} explains how to build and verify them.
